\chapter{Ход работы}
\label{ch:chap2}

\section{Реализация многоканальности}

Для реализации многоканальности используются блоки Resource Pool, Resource Acquirer, 
Resource Releaser из библиотеки SimEvent.

\subsection{Resource pool}

Специальный блок, который воспринимает блоки Simulink как ресурсы. Самое важное в этом блоке то, что он может считать, что ресурсов
у него несколько, хотя блок фактически только один. Благодря этому блоку, мы можем задать в системе N штук ресурсов с именем server (например).
Эти ресурсы можно занимать и осовобождать. Сам Resource pool не освобождает и не занимает ресусры, а только ведет их учет (сколько занято,
сколько свободно).
 
\begin{figure}[H]
    \centering
    \includegraphics[width=1.0\textwidth]{resource_pool.png}
    \caption{Resource poll}
\end{figure}

В окне настроек можно задать Resource amount, т.е максимальное кол-во ресурсов.


\subsection{Resource Acquirer}

Блок, который занимает ресурс. В нашем случае при поступлении в этот блок заявки, он выделит на ее обработку один ресурс (сервер).


\begin{figure}[H]
    \centering
    \includegraphics[width=1.0\textwidth]{res_ac.png}
    \caption{Resource Acquirer}
\end{figure}

В окне настройки необходимо указать в поле Name имя блока, который мы хотим рассматривать как ресурс (у меня это entity server с именем Server)

\subsection{Resource Releaser}

Блок, который освобождает ресурс. Задача этого блока в том, чтобы сказать Resource pool, что ресурс выполнил задачу по обработке
заявки и он готов к работе. \\

Дополнительных настроек блок не требует.

\subsection{Создание подсистемы}

Чтобы дать пользователю возможность регулировать кол-во серверов и вероятность отказа, необходимо создать подсистему, в которую потом
пользователь сможет передавать параметры. Для этого необходимо выделить блоки, которые нужно поместить в подсистему, и в появившемся
окне нажать на "Create subsystem".

\begin{figure}[H]
    \centering
    \includegraphics[width=1.0\textwidth]{subsys.png}
    \caption{Подсистема}
\end{figure}

\subsection{Настройка масок подсистемы}

Под маской в Simulink подразумеваются передаваемые в подсистему параметры. Для настройки маски нужно нажать ПКМ на подсистему и выбрать
пункт "Maks". В открывшемся окне перейти в PARAMETERS\&DIALOG перетащить блоки "edit" (слева) в окно PARAMETERS\&DIALOG. Каждый блок
"edit" - переменная, которую может передавать в подсистему пользователь. Этот блок нужно настроить. Для этого в свойствах блока (справа)
необходимо указать свойства Name, value, prompt. Name - имя, по которому будет доступна переменная в подсистеме, value - текущее значение
переменной, prompt - описание, которое видит пользователь.

\begin{figure}[H]
    \centering
    \includegraphics[width=1.0\textwidth]{mask.png}
    \caption{Маска подсистемы}
\end{figure}

В результате у пользователя должна появиться возможность настраивать работу подсистемы

\begin{figure}[H]
    \centering
    \includegraphics[width=1.0\textwidth]{res.png}
    \caption{Настройка подсистемы}
\end{figure}

\subsection{Вывод данных из подсистемы}

Для вывода данных из подсистемы нужен блок "out1". После добавления этого блока на выходе подсистемы появится выход.

\section{Реализация резервных серверов}

Добавим еще одну такую же подсистему, но на нее будут поступать заявки не напрямую из очереди, а только в том случае, если заявка
была отброшена основной подсистемой. Если заявка не проходит резервную подсистему, то по обратной связи заявки поступают обратно в очередь.

\section{Общий вид системы M/M/c}

\begin{figure}[H]
    \centering
    \includegraphics[width=1.0\textwidth]{mmc.png}
    \caption{Общий вид системы}
\end{figure}

\section{Тесты}

\subsection{Без резервирования}

\subsubsection{Один сервер}

\begin{figure}[H]
    \centering
    \includegraphics[width=1.0\textwidth]{1_serv.png}
    \caption{Тест с одним сервером}
\end{figure}

Рассмотрим влияние кол-ва серверов на очередь СМО. Для это запустим симуляцию в кол-вом серверов (n) равным 1, вероятностью
отказа (Potk) равной 0.5

\subsubsection{Пять серверов}

Запустим симуляцию повторно, но уже с n = 5

\begin{figure}[H]
    \centering
    \includegraphics[width=1.0\textwidth]{5_serv.png}
    \caption{Тест с пятью серверами}
\end{figure}


Можем заметить, что средняя длина очереди и среднее время ожидания заявок в очреди стали близки к нулю. Из этого можно сделать
вывод о том, что кол-во потоков обработки заявок в СМО напрямую влияет на скорость обслуживания заявок.

\subsection{С резервированием}

Для демонстрации важности резервирования серверов поставим Potk главной подсистемы равной 0.9. На резервной системе вероятность
отказа составляет 0.2

\subsubsection{1 резервный сервер}

Предположим, что у нас 1 главный сервер и 1 резервный. Запустим симуляцию

\begin{figure}[H]
    \centering
    \includegraphics[width=1.0\textwidth]{1res.png}
    \caption{Тест с 1 резервным сервером}
\end{figure}

\subsubsection{10 резервных серверов}

Увеличим кол-во резервных серверов и запустим симуляцию

\begin{figure}[H]
    \centering
    \includegraphics[width=1.0\textwidth]{10res.png}
    \caption{Тест с 10 резервными серверами}
\end{figure}

Можем заметить, что средняя длина очереди и среднее время ожидания заявки в очереди стало меньше. Это говорит о том, что резервирование
серверов способствует улучшению показателей СМО в случае большого числа отказов главного сервера.


\endinput
